\paragraph{去极化、原始化与 dyadic 零点输运的三层分解}
前述二进尺度递归已经把 \(\Xi\)-函数的截断误差与简单零点稳定性压缩为显式超指数预算。
若再把离散 Abel--Weil 生成函数上的幂覆盖去极化、M\"obius 多尺度中和、Hardy 通道能量分解与 completion-side 的 dyadic 余项相位律并置，则可把“值层逼近”与“零点敏感输运”之间的缺口压缩为一个单独的相位寄存器问题。

\begin{theorem}[去极化与 M\"obius 原始化的互补正交性]\label{thm:conclusion-xi-depolarization-mobius-orthogonality}
\leanverified{paper\_conclusion\_xi\_depolarization\_mobius\_orthogonality}
固定 \(b>1\)。
令
\[
\mathcal A_b^+(r):=\mathcal A_b(r)-\mathcal A_b(0),
\qquad
\Delta_m\mathcal A_b(r):=\mathcal A_b(r)-m\,\mathcal A_{b^m}(r^m),
\qquad
\mathfrak M_b(r):=\sum_{d\ge 1}\mu(d)\,\mathcal A_{b^d}^+(r^d),
\]
并记
\[
K_m(u):=\frac{1}{1-u}-\frac{m}{1-u^m}.
\]
则对任意整数 \(m\ge 2\)：
\begin{enumerate}
  \item 去极化恒等式
  \[
  \Delta_m\mathcal A_b(r)
  =
  -\sum_{\rho}\frac{1}{\rho}\,K_m\!\bigl(r\,b^{\rho-\frac12}\bigr)
  +\bigl(H_b(r)-mH_{b^m}(r^m)\bigr)
  \]
  成立，且
  \[
  K_m(1)=-\frac{m-1}{2}.
  \]
  因而原先位于 \(u=1\) 的极点并不保留在径向位置，而是被推送到
  \[
  u^m=1,\qquad u\neq 1
  \]
  的非平凡根单位轨道。

  \item 在 RH 条件下，M\"obius 原始化满足
  \[
  \mathfrak M_b(r)
  =
  -r\sum_{\rho}\frac{b^{\rho-\frac12}}{\rho}
  +\mathfrak H_b(r),
  \]
  即全部高倍几何级数尾项被完全中和，只余第一层 primitive 谱壳。
\end{enumerate}
特别地，去极化与 M\"obius 原始化处理的是两类严格正交的冗余：前者作用于 cyclotomic 根单位轨道，后者作用于几何倍层。
\end{theorem}
\begin{proof}
第一条正是定理 \ref{thm:abel-depolarization-root-unity-rotation} 的内容。
第二条正是定理 \ref{thm:abel-mobius-linearization} 的线性化结论。
两式并列可见：\(\Delta_m\) 并不消去 primitive 一次谱项，而 \(\mathfrak M_b\) 也不处理 \(u^m=1\) 的根单位旋转，因此二者分别剥离的是 cyclotomic 轨道与几何倍层这两类互补方向。证毕。
\end{proof}

\begin{theorem}[方向均匀化与总体 Hardy 缺陷的精确分裂]\label{thm:conclusion-xi-direction-total-defect-splitting}
\leanverified{paper\_conclusion\_xi\_direction\_total\_defect\_splitting}
对 \(\delta>0\) 定义加权 Hardy 生成函数
\[
\mathcal A_{b,\delta}(r):=\sum_{n\ge 0}E_b(n)b^{-\delta n}r^n.
\]
固定 \(m\ge 2\)。
则存在如下严格分裂：
\begin{enumerate}
  \item 对每个非平凡加性特征 \(\chi:\ZZ/m\ZZ\to\mathbb T\)，归一化字符相关
  \[
  \mathcal F_\chi(\delta):=
  \frac{\sum_{n\ge 0}|E_b(n)|^2b^{-2\delta n}\chi(n)}
  {\sum_{n\ge 0}|E_b(n)|^2b^{-2\delta n}}
  \]
  在 \(\delta\downarrow 0\) 时趋于 \(0\)；等价地，polyphase 能量分布趋于
  \[
  \left(\frac1m,\dots,\frac1m\right).
  \]

  \item 然而与此同时，去极化 Hardy 缺陷满足精确极限
  \[
  \lim_{\delta\downarrow 0}
  \frac{\|\Delta_m\mathcal A_{b,\delta}\|_{H^2(\mathbb D)}^2}
  {\|\mathcal A_{b,\delta}\|_{H^2(\mathbb D)}^2}
  =
  m-1.
  \]
\end{enumerate}
因此，每个固定的非平凡 cyclotomic 方向都在极限中消失，但总体 Hardy 缺陷并不消失，而是以精确比例 \(m-1\) 存活。
\end{theorem}
\begin{proof}
第一条即该节 polyphase 能量均匀化定理在角色侧的表述；由命题 \ref{prop:abel-channel-equipartition-character}，它等价于全部非平凡字符系数同时趋于零。
第二条则正是定理 \ref{thm:abel-depolarization-hardy-energy-blowup} 的倍率极限。
二者合并便得到“方向消失而总缺陷存活”的严格分裂。证毕。
\end{proof}

\begin{corollary}[有限方向账本的绝对失明]\label{cor:conclusion-xi-finite-direction-ledger-blindness}
\leanverified{paper\_conclusion\_xi\_finite\_direction\_ledger\_blindness}
任何仅通过有限个非平凡 residue-class 特征通道来审计 \(\mathcal A_{b,\delta}\) 的协议，在 \(\delta\downarrow 0\) 时都必然把对象读成“渐近均匀”；然而任何读取整体 \(H^2\)-缺陷的协议都会检测到一个精确 surviving defect，其相对规模恒为 \(m-1\)。
因此，有限方向账本不能替代全局 Hardy 缺陷账本。
\end{corollary}
\begin{proof}
由定理 \ref{thm:conclusion-xi-direction-total-defect-splitting}，每个固定非平凡特征的归一化相关都趋于 \(0\)，故任意有限方向采样最终都失去分辨力。
而同一定理同时给出整体 Hardy 缺陷比值的极限恰为 \(m-1\)。
于是方向信息全部洗平并不意味着对象级缺陷消失。证毕。
\end{proof}

\begin{theorem}[Riemann--Siegel 阶数链的 Cauchy 刚性]\label{thm:conclusion-xi-order-chain-cauchy-rigidity}
\leanverified{paper\_conclusion\_xi\_order\_chain\_cauchy\_rigidity}
设 \(\{t_m\}_{m\ge m_0}\) 为一条 admissible 的局部 Riemann--Siegel 截断零点链，其中 \(t_m\) 是 \(m\) 阶局部截断 \(Z^{(m)}\) 的一个简单零点，并假设该零点链沿阶数持续满足正文 13.4238--13.4240 所要求的局部可输运条件。
则存在一步传输界
\[
|t_{m+1}-t_m|
\ll
\frac{t_m^{-(m+1)/2}}{\log t_m}.
\]
特别地，只要 admissibility 沿阶数持续成立，则 \(\{t_m\}\) 在阶数方向绝对可和，从而为 Cauchy 序列；其极限由任意 admissible 阶数链唯一决定。
\end{theorem}
\begin{proof}
一步传输估计正是正文 13.4238--13.4240 的局部零点输运结论。
对固定大高度 \(t_\ast\)，右端在 \(m\) 上以 \(t_\ast^{-1/2}\) 为公比形成几何级数，因而
\[
\sum_{m\ge m_0}|t_{m+1}-t_m|<\infty.
\]
于是 \(\{t_m\}\) 为 Cauchy 序列。
同理，任意两条 admissible 阶数链之间的差异也受同一尾和控制，故极限唯一。证毕。
\end{proof}

\begin{theorem}[二进递归下简单零点的首项漂移律]\label{thm:conclusion-xi-dyadic-zero-phase-drift}
设
\[
\rho_0=\frac12+\mathrm{i}\tau
\]
是 \(\xi\) 的一个简单零点，\(\xi_K\) 为完成核的 dyadic 截断，\(\rho_K\) 为定理 \ref{thm:xi-completed-zeta-simple-zero-stability} 在 \(\rho_0\) 邻域内给出的唯一简单零点。
则有一阶漂移展开
\[
\rho_K-\rho_0
=
\frac{R_K(\rho_0)}{\xi'(\rho_0)}+o_{\rho_0}\!\bigl(R_K(\rho_0)\bigr),
\]
其中
\[
R_K\!\left(\frac12+\mathrm{i}\tau\right)
=
\frac{2}{\pi}\,2^{-3K/4}e^{-\pi2^K}
\cos\!\left(\frac{\tau K\log2}{2}\right)
+O_\tau\!\left(2^{-7K/4}e^{-\pi2^K}\right)
+O\!\left(2^{-3K/4}e^{-4\pi2^K}\right).
\]
因而
\[
\rho_K-\rho_0
=
\frac{2}{\pi\,\xi'(\rho_0)}\,2^{-3K/4}e^{-\pi2^K}
\cos\!\left(\frac{\tau K\log2}{2}\right)
+o_{\rho_0}\!\left(2^{-3K/4}e^{-\pi2^K}\right).
\]
因此 dyadic 零点输运的首项并非仅由尾项振幅控制，而由明确的二进相位
\[
\frac{\tau K\log2}{2}
\]
锁定。
\end{theorem}
\begin{proof}
由定理 \ref{thm:xi-completed-zeta-simple-zero-stability}，\(\rho_K\to\rho_0\) 且 \(\xi_K(\rho_K)=0\)。
又 \(\xi=\xi_K+R_K\) 且 \(\xi(\rho_0)=0\)，故
\[
0=\xi_K(\rho_K)=\xi(\rho_K)-R_K(\rho_K).
\]
在 \(\rho_0\) 处对 \(\xi\) 作 Taylor 展开，并利用 \(\xi'(\rho_0)\neq0\)，得到
\[
\rho_K-\rho_0=\frac{R_K(\rho_0)}{\xi'(\rho_0)}+o_{\rho_0}\!\bigl(R_K(\rho_0)\bigr).
\]
再代入定理 \ref{thm:xi-completed-zeta-remainder-phase-law} 的余项主项展开即得。证毕。
\end{proof}

\begin{corollary}[零点稳定的相位统治]\label{cor:conclusion-xi-dyadic-phase-governed-stability}
在定理 \ref{thm:conclusion-xi-dyadic-zero-phase-drift} 的设定下：
\begin{enumerate}
  \item 若
  \[
  \frac{\tau\log2}{2\pi}\notin\QQ,
  \]
  则主相位
  \[
  \cos\!\left(\frac{\tau K\log2}{2}\right)
  \]
  沿 \(K\in\NN\) 在 \([-1,1]\) 中稠密，因而其符号改变无穷多次。

  \item 若
  \[
  \frac{\tau\log2}{2\pi}\in\QQ,
  \]
  则首项漂移呈现有限周期相位律。
\end{enumerate}
因此超指数稳定只控制振幅，不控制方向；零点漂移的主导行为 generically 由相位而非单调性统治。
\end{corollary}
\begin{proof}
由定理 \ref{thm:conclusion-xi-dyadic-zero-phase-drift}，漂移的首项完全由
\(
\cos(\tau K\log2/2)
\)
控制。
若 \(\tau\log2/(2\pi)\) 无理，则由 Kronecker 稠密定理可得相位序列在单位圆上稠密，故余弦值在 \([-1,1]\) 中稠密并无穷次变号。
若该比值有理，则相位序列周期，结论随之成立。证毕。
\end{proof}

\begin{proposition}[自然边界下的有限阶递归证书失效]\label{prop:conclusion-xi-natural-boundary-no-finite-recursive-certificate}
\leanverified{paper\_conclusion\_xi\_natural\_boundary\_no\_finite\_recursive\_certificate}
在 RH 及
\[
\mathcal R_b=\left\{e^{-\mathrm{i}\gamma\log b}:\ \zeta\!\left(\frac12+\mathrm{i}\gamma\right)=0\right\}
\]
在单位圆上稠密的条件下，\(\mathcal A_b(r)\) 以单位圆为自然边界，并且误差序列 \(\{E_b(n)\}\) 不存在任何
\[
\text{“有限阶常系数递归}+\text{严格指数误差”}
\]
的闭合模型。
因此，任何试图以有限阶线性状态机忠实压缩该误差层的证书范式都不可能成功。
\end{proposition}
\begin{proof}
自然边界结论正是定理 \ref{thm:xi-abel-rh-dense-poles-natural-boundary} 的内容。
而推论 \ref{cor:xi-abel-natural-boundary-no-finite-order-exp-lock} 进一步排除了任意固定阶常系数递归配合严格指数误差的闭包。
两者合并便得到所述失效原理。证毕。
\end{proof}

\begin{theorem}[二进相位寄存器下界]\label{thm:conclusion-xi-dyadic-phase-register-lower-bound}
\leanverified{paper\_conclusion\_xi\_dyadic\_phase\_register\_lower\_bound}
考虑如下证书模型类：其内部状态由有限维常系数线性递推更新；其对零点的预测只允许附加严格指数衰减的残差；并且不显式携带随 \(K\) 演化的外置相位寄存器。
则该模型类不可能忠实实现定理 \ref{thm:conclusion-xi-dyadic-zero-phase-drift} 的 dyadic zero transport。

更精确地说，若某模型 \(\mathcal M\) 在每个简单零点附近都能再现 \(\rho_K-\rho_0\) 的首项行为，则 \(\mathcal M\) 至少必须外置一个相位寄存器
\[
\Phi_K(\tau)\equiv \frac{\tau K\log2}{2}\pmod{2\pi}.
\]
\end{theorem}
\begin{proof}
若不存在这样的外置相位寄存器，则模型输出只能由有限阶常系数状态与严格指数残差决定。
但命题 \ref{prop:conclusion-xi-natural-boundary-no-finite-recursive-certificate} 已排除 projective error layer 上的这类有限阶闭包。
另一方面，定理 \ref{thm:conclusion-xi-dyadic-zero-phase-drift} 的首项含有不可消去的相位因子
\[
\cos\!\left(\frac{\tau K\log2}{2}\right),
\]
它既不由振幅 \(2^{-3K/4}e^{-\pi2^K}\) 决定，也不由有限维常系数状态自动生成。
故若不外置 \(\Phi_K(\tau)\)，则模型无法忠实再现不同 \(\tau\) 对应的 dyadic 相位输运。证毕。
\end{proof}

\begin{conjecture}[零点敏感 \texorpdfstring{$\pi$}{pi}-通道的三分解原则]\label{conj:conclusion-xi-zero-sensitive-pi-channel-trisplitting}
任一不仅在数值上逼近 \(\pi\)，而且能够忠实承载 completion-side 零点输运的 \(\pi\)-通道，应至少分解为三层：
\[
\mathfrak V_{\pi}
\quad\text{（标量值层，负责快速数值逼近）},
\]
\[
\mathfrak P_{\mathrm{prim}}
\quad\text{（primitive 谱层，负责 M\"obius 原始化）},
\]
\[
\mathfrak R_{\mathrm{dyad}}
\quad\text{（二进相位寄存器，负责 }K\mapsto \tau K\log2/2\text{ 的相位输运）}.
\]
缺失第三层的公式，即便在数值意义下极其高效，也至多是值层算法，而不是零点敏感的 completion 证书。
\end{conjecture}
